\documentclass[11pt]{article}
\usepackage[a4paper,margin=22mm]{geometry}
\usepackage{amsmath,amssymb,amsthm,mathtools,bm}
\usepackage{slashed}
\usepackage{booktabs,longtable,array}
\usepackage{enumitem}
\usepackage{xcolor}
\usepackage{microtype}
\usepackage[hidelinks]{hyperref}

\setlength{\emergencystretch}{3em}
\newtheorem{theorem}{Theorem}[section]
\newtheorem{proposition}[theorem]{Proposition}
\newtheorem{lemma}[theorem]{Lemma}
\newtheorem{corollary}[theorem]{Corollary}
\theoremstyle{definition}
\newtheorem{definition}[theorem]{Definition}
\newtheorem{gate}[theorem]{Fail-closed gate}
\theoremstyle{remark}
\newtheorem{remark}[theorem]{Remark}

\newcommand{\dd}{\mathrm d}
\newcommand{\ii}{\mathrm i}
\newcommand{\ZZ}{\mathbb Z}
\newcommand{\RR}{\mathbb R}
\newcommand{\CC}{\mathbb C}
\newcommand{\Tr}{\operatorname{Tr}}
\newcommand{\tr}{\operatorname{tr}}
\newcommand{\ind}{\operatorname{ind}}
\newcommand{\SF}{\operatorname{SF}}
\newcommand{\Pf}{\operatorname{Pf}}
\newcommand{\sgn}{\operatorname{sgn}}
\newcommand{\Sym}{\operatorname{Sym}}
\newcommand{\status}[1]{\textcolor{blue!60!black}{\textsf{#1}}}
\newcolumntype{L}[1]{>{\raggedright\arraybackslash}p{#1}}

\title{AP-E6: The Simply Connected \(Sp(4)\) Candidate\\
Partial Euclidean Controls, Gauge Bordism, Target-Space Eta No-Go,\\
and Heavy-Threshold Audit}
\author{Route F independent research note}
\date{16 July 2026}

\begin{document}
\maketitle

\begin{abstract}
We complete a sharply delimited audit of the proposed simply connected
deep group
\(Sp(4)_{\rm phys}=USp(4)=Spin(5)=Sp(2)_{\rm math}\).  First, we give a
Euclidean action ansatz and verify two limited fermionic controls for two
Dirac \(\bm5\)'s: gamma-five reality/positivity and three Pauli--Villars
moment identities.  This is \emph{not} a complete Euclidean regulator.
In particular, the regulator-field statistics and full mass interpolation,
APS boundary data, finite local-counterterm scheme, and regulated
gauge--scalar--ghost measure have not been specified.

Second, an explicit Atiyah--Hirzebruch spectral-sequence computation gives
\[
 \boxed{\Omega^{\mathrm{Spin}}_5(BSp(2))\simeq\ZZ_2.}
\]
Its generator is an embedded unit \(Sp(1)\) instanton on \(S^4\), times the
periodic spin circle, equivalently the mapping torus of the nontrivial
element of \(\pi_4(Sp(2))\).  Restriction to this \(Sp(1)\) gives mod-two
indices \(\nu_{\bm4}=1\), \(\nu_{\bm5}=0\), and
\(\nu_{\bm{10}}=0\).  Hence the displayed two Dirac \(\bm5\)'s define the
trivial character of the nonzero bordism class.  The dynamical
\(Sp(4)\) gauge-bordism anomaly gate is therefore closed for this field
list.

Third, this result does \emph{not} select the two target-space signs on
\(G_3=S^1_{\rm R}\times S^3\) and
\(G_2=T^2_{\rm RR}\times S^2\).  Those signs are characters of
\(\widetilde\Omega^{\mathrm{Spin}}_4(S^3\times S^2)\), and require a mass
map from the sigma-model target into the regulated Fredholm family.  The
current \(Sp(4)\) action specifies neither the composite pion lift nor the
full \(s\)-dependent Dirac--Yukawa/PV family.  We prove underdetermination
by enumerating the independent torsion invertible sectors/counterterms
that flip either character while leaving the gauge character and de Rham
curvature fixed.  This enumeration is not a microscopic realization of
those sectors.

Finally, threshold bookkeeping produces a stronger obstruction.  The
\(\bm5\) branches as
\(\bm2_{+1}\oplus\bm2_{-1}\oplus\bm1_0\).  Uniform flavour action on the
complete irrep gives
\(2(+1)+2(-1)+0=0\).  Under the additional hypothesis that the chiral
flavour symmetry acts uniformly and exactly on the complete \(\bm5\), the
light/heavy charge traces are conditionally \(+2\) and \(-2\).  This is an
algebraic obstruction, not an evaluated heavy determinant.
The physical \(B=+1\) infrared orientation \(k_{\rm IR}=+2\) is derived,
but it can survive only if the chiral flavour symmetry is emergent below
the breaking threshold or if new UV topological data are added.  Neither
option is yet a closed first-principles completion.  No Route-E portal is
authorized.
\end{abstract}

\tableofcontents

\section{Scope, notation, and exact claim boundary}

Throughout this note ``\(Sp(4)\)'' means the physics group with a
four-dimensional defining representation,
\begin{equation}
 G=Sp(4)_{\rm phys}=USp(4)=Spin(5)=Sp(2)_{\rm math}.
 \label{eq:notation}
\end{equation}
This convention is essential when reading bordism tables.  The candidate
field content is
\begin{equation}
 \Psi_i\in\bm5\ \text{Dirac},\quad i=1,2;
 \qquad
 \Phi_I\in\bm4\ \text{complex},\quad I=1,2;
 \qquad
 \Sigma\in\bm{10}\ \text{real}.
 \label{eq:fields}
\end{equation}
The copy index \(I\) carries the desired global \(SU(2)_\phi\); it is not a
gauge weight.  The standard branch through
\(Spin(4)=SU(2)_c\times SU(2)_H\) is
\begin{align}
 \bm4&\longrightarrow(\bm2,\bm1)\oplus(\bm1,\bm2),
 \label{eq:four-spin4}\\
 \bm5&\longrightarrow(\bm2,\bm2)\oplus(\bm1,\bm1),
 \label{eq:five-spin4}\\
 \bm{10}&\longrightarrow(\bm3,\bm1)\oplus(\bm1,\bm3)
                       \oplus(\bm2,\bm2).
 \label{eq:ten-spin4}
\end{align}
After \(SU(2)_H\to U(1)_H\), with doublet weights normalized to
\(q=\pm1\),
\begin{align}
 \bm4&\to\bm2_0\oplus\bm1_{+1}\oplus\bm1_{-1},
 \label{eq:four-u1}\\
 \bm5&\to\bm2_{+1}\oplus\bm2_{-1}\oplus\bm1_0,
 \label{eq:five-u1}\\
 \bm{10}&\to\bm1_{+2}\oplus\bm2_{+1}\oplus\bm3_0\oplus\bm1_0
                 \oplus\bm2_{-1}\oplus\bm1_{-2}.
 \label{eq:ten-u1}
\end{align}
These identities can be obtained from
\(\bm5=\wedge^2\bm4-\bm1\) and
\(\bm{10}=\Sym^2\bm4\) \cite{Slansky1981}.

Four logically different questions will not be conflated:
\begin{enumerate}[label=(\roman*)]
\item whether the Euclidean \(Sp(4)\) path integral has been specified;
\item whether its \emph{dynamical gauge} determinant line has a global
  anomaly, classified here by \(\Omega^{\rm Spin}_5(BSp(2))\);
\item which spin-torsion characters the induced \emph{target-space}
  functional assigns to \(G_3\) and \(G_2\);
\item whether the heavy threshold preserves the desired integer mixed-WZW
  coefficient and the \(k=+2\) orientation.
\end{enumerate}
Only (ii) will close.  Question (i) has only partial fermion controls,
(iii) is proved underdetermined, and (iv) has a conditional uniform-flavour
trace obstruction but no APS determinant-ratio computation.  This separation follows the
APS/Dai--Freed treatment of determinant phases
\cite{AtiyahPatodiSinger1975,DaiFreed1994,WittenYonekura2019}.

\section{Partial Euclidean fermion-regulator controls}

\subsection{Bundles, hermiticity, and the bosonic action}

Let \(M_4\) be a closed oriented Riemannian spin manifold and
\(P\to M_4\) a principal \(G\)-bundle.  Write
\(E_R=P\times_G R\), choose Hermitian Euclidean gamma matrices, and use an
anti-Hermitian gauge connection \(A\).  If \(T_R^a\) are anti-Hermitian
generators, define the Hermitian adjoint-Higgs endomorphism
\begin{equation}
 \widehat\Sigma_R:=-\ii\Sigma^aT_R^a,
 \qquad \widehat\Sigma_R^\dagger=\widehat\Sigma_R.
 \label{eq:sigma-hat}
\end{equation}
The gauge-fixed bosonic action may be chosen as
\begin{align}
 S_{E,{\rm bos}}={}&\int_{M_4}\!\dd^4x\sqrt g\,
 \mathcal L_{\rm bos}+S_{\rm gf}+S_{\rm gh},
 \label{eq:bosonic-action}\\
 \mathcal L_{\rm bos}={}&
 \frac1{2g_G^2}\tr(F_{\mu\nu}^\dagger F^{\mu\nu})
 +\frac12(D_\mu\Sigma)^a(D^\mu\Sigma)^a
 \notag\\
 &
 +\sum_{I=1}^2(D_\mu\Phi_I)^\dagger D^\mu\Phi_I
 +V(\Sigma,\Phi),
 \label{eq:bosonic-density}\\
 V={}&\frac{\lambda_\Sigma}{4}
   \bigl(\Sigma^a\Sigma^a-v_\Sigma^2\bigr)^2
 +m_\Phi^2\sum_I\Phi_I^\dagger\Phi_I
 +\kappa\sum_I\Phi_I^\dagger\widehat\Sigma_{\bm4}\Phi_I
 \notag\\
 &+\lambda_1\Bigl(\sum_I\Phi_I^\dagger\Phi_I\Bigr)^2
 +\lambda_2\sum_{I,J}|\Phi_I^\dagger\Phi_J|^2.
 \label{eq:potential}
\end{align}
Here \(\lambda_\Sigma,\lambda_1>0\) and
\(\lambda_1+\lambda_2>0\) give a simple bounded quartic domain.  The
copy-symmetric coefficients preserve \(SU(2)_\phi\).  A background
\(R_\xi\) gauge fixing and its Faddeev--Popov determinant define
\(S_{\rm gf}+S_{\rm gh}\) formally; no regulator for the
gauge--scalar--ghost complex is supplied here.  This omission does not
affect the separate bordism computation below, but it keeps the full
Euclidean-regulator gate open.

Choose an adjoint vacuum
\begin{equation}
 \langle\widehat\Sigma\rangle=V H,
 \qquad H=-2\ii T_H^3,
 \label{eq:adjoint-vev}
\end{equation}
where \(T_H^3\) is anti-Hermitian.  Thus \(H\) is Hermitian and has the
real weights displayed in
Eqs.~\eqref{eq:four-u1}--\eqref{eq:ten-u1}.  Because
\(Spin(4)\hookrightarrow Spin(5)\) is injective, its central element
\((-\bm1_c,-\bm1_H)\) maps to the nontrivial centre of \(Spin(5)\), not to
the identity.  The unbroken group is consequently the direct product
\(SU(2)_c\times U(1)_H\), rather than its diagonal \(\ZZ_2\) quotient.

\subsection{Dirac--Yukawa operator and threshold masses}

For each flavour let
\begin{equation}
 \mathcal K_{A,\Sigma}
 :=\gamma^\mu(\nabla_\mu+A_\mu^{(\bm5)})
 +m\bm1_{\bm5}+y\widehat\Sigma_{\bm5}.
 \label{eq:dirac-yukawa}
\end{equation}
The fermion action is
\begin{equation}
 S_{E,{\rm f}}=\sum_{i=1}^2
 \int_{M_4}\!\dd^4x\sqrt g\,\overline\Psi_i\mathcal K_{A,\Sigma}\Psi_i.
 \label{eq:fermion-action}
\end{equation}
On the three branches of \(\bm5\),
\begin{equation}
 M_+=m+yV,
 \qquad M_-=m-yV,
 \qquad M_0=m.
 \label{eq:masses-general}
\end{equation}
At the tree-level light-point \(m=-yV\), and with \(yV>0\),
\begin{equation}
 (M_+,M_-,M_0)=(0,-2yV,-yV).
 \label{eq:masses-tuned}
\end{equation}
For determinant statements we keep a positive infrared mass
\(M_+=\mu>0\) and take \(\mu\downarrow0\) only after all phases and zero
modes have been recorded.

The operator obeys
\begin{equation}
 \mathcal K_{A,\Sigma}^\dagger
 =\gamma_5\mathcal K_{A,\Sigma}\gamma_5.
 \label{eq:gamma5-hermiticity}
\end{equation}
Thus its nonreal eigenvalues occur in complex-conjugate pairs and
\(\det\mathcal K\in\RR\).  For the two identical flavours,
\begin{equation}
 Z_{\Psi,\mu}=\bigl(\det\mathcal K_{A,\Sigma;\mu}\bigr)^2\ge0
 \label{eq:two-flavour-positive}
\end{equation}
away from the declared zero locus.  This is a two-flavour positivity
statement, not a claim that an arbitrary chiral mass family has positive
determinant.

\subsection{Pauli--Villars amplitude and eta phase}

It is useful to separate the ultraviolet amplitude subtraction from the
global phase.  Put
\begin{equation}
 c_a=(1,-3,3,-1),\qquad
 \Lambda_a^2=a\Lambda^2,\qquad a=0,1,2,3,
 \label{eq:pv-data}
\end{equation}
The finite-difference identities
\begin{equation}
 \sum_ac_a=\sum_ac_aa=\sum_ac_aa^2=0
 \label{eq:pv-moments}
\end{equation}
remove the first three heat-kernel power moments in
\begin{equation}
 \log|\det\mathcal K|_{\rm PV}
 =\frac12\sum_{a=0}^3c_a
   \log\det\bigl(\mathcal K^\dagger\mathcal K+\Lambda_a^2\bigr).
 \label{eq:pv-amplitude}
\end{equation}
The remaining logarithmic divergences would have to be removed by a
declared finite gauge-invariant local-counterterm scheme.  Such a scheme is
not supplied in this note.

For the phase, a completed construction would choose a five-dimensional
interpolation \(Y_5\), self-adjoint operators \(D_{Y,R,a}\), compatible APS
boundary conditions, regulator statistics, and PV reference signs
\(\sigma_a\in\{\pm1\}\).  Its torsion datum would have the schematic form
\begin{equation}
 \tau_R(Y_5;\sigma)
 =\exp\!\left[2\pi\ii\sum_ac_a\sigma_a
 \overline\eta(D_{Y,R,a})\right],
 \qquad
 \overline\eta=\frac{\eta+\dim\ker D}{2}.
 \label{eq:eta-regulator}
\end{equation}
Once all boundary and interpolation data are supplied, APS gluing can make
the torsion part a bordism character.  Here
Eq.~\eqref{eq:eta-regulator} is only a template: neither the full
\(\mathcal K(t),\Lambda_a(t)\) mass interpolation nor the APS boundary
operator/domain, local counterterms, and gauge--scalar--ghost regulator are
defined.

\begin{proposition}[Partial fermion-regulator controls]
Equations~\eqref{eq:gamma5-hermiticity}--\eqref{eq:pv-moments} establish
gamma-five reality, two-flavour positivity away from the declared zero
locus, and three PV moment identities.  They do not define the complete
functional measure or its phase.  Consequently
\begin{equation}
 \texttt{fermion\_pv\_gamma5\_controls\_verified=true},\qquad
 \texttt{sp4\_euclidean\_regulator\_complete=false}.
 \label{eq:regulator-partial-gate}
\end{equation}
\end{proposition}

\section{The exact gauge-bordism group}

\subsection{AHSS computation}

The integral cohomology ring is
\begin{equation}
 H^*(BSp(2);\ZZ)=\ZZ[q_1,q_2],
 \qquad |q_1|=4,\quad |q_2|=8.
 \label{eq:bsp-cohomology}
\end{equation}
Consequently, in degrees below eight,
\begin{equation}
 H_p(BSp(2);\ZZ)=
 \begin{cases}
 \ZZ,&p=0,4,\\
 0,&0<p<8,\ p\ne4,
 \end{cases}
 \label{eq:bsp-homology}
\end{equation}
and \(H_4(BSp(2);\ZZ_2)=\ZZ_2\).  The homological AHSS is
\begin{equation}
 E^2_{p,q}=H_p(BSp(2);\Omega_q^{\rm Spin})
 \Longrightarrow\Omega_{p+q}^{\rm Spin}(BSp(2)),
 \qquad
 d_r:E^r_{p,q}\to E^r_{p-r,q+r-1}.
 \label{eq:ahss}
\end{equation}
Use
\begin{equation}
 \Omega_0^{\rm Spin}=\ZZ,\quad
 \Omega_1^{\rm Spin}=\ZZ_2,\quad
 \Omega_2^{\rm Spin}=\ZZ_2,\quad
 \Omega_3^{\rm Spin}=0,\quad
 \Omega_4^{\rm Spin}=\ZZ,\quad
 \Omega_5^{\rm Spin}=0.
 \label{eq:spin-coefficients}
\end{equation}
On the total-degree-five diagonal the only nonzero term is
\begin{equation}
 E^2_{4,1}=H_4(BSp(2);\ZZ_2)=\ZZ_2.
 \label{eq:unique-ahss-term}
\end{equation}
The outgoing maps are
\begin{align}
 d_2:E^2_{4,1}&\longrightarrow E^2_{2,2}=0,
 \notag\\
 d_3:E^3_{4,1}&\longrightarrow E^3_{1,3}=0,
 \notag\\
 d_4:E^4_{4,1}&\longrightarrow E^4_{0,4}=\ZZ.
 \label{eq:outgoing}
\end{align}
The last homomorphism vanishes because there is no nonzero group
homomorphism \(\ZZ_2\to\ZZ\).  There is no incoming \(d_2\), since
\(E^2_{6,0}=H_6(BSp(2);\ZZ)=0\), and for \(r\ge3\) an incoming
differential would start at negative coefficient degree.  Hence the class
survives to \(E^\infty\), and there is no extension problem because it is
the only filtration quotient.

\begin{theorem}[Low-degree \(Sp(4)\) gauge bordism]
For the simply connected group in Eq.~\eqref{eq:notation},
\begin{equation}
 \boxed{\Omega_5^{\rm Spin}(BSp(2))\simeq\ZZ_2.}
 \label{eq:omega5-result}
\end{equation}
\end{theorem}
This agrees with the direct \(Sp(n)\) analysis of Saito and Tachikawa
\cite{SaitoTachikawa2025}; recent instanton zero-mode analysis is also
given by Jia and Yi \cite{JiaYi2024}.

\subsection{Geometric generator and representation character}

Let \(P_1\to S^4\) be the unit instanton bundle for an embedded
\(Sp(1)=SU(2)_c\subset Sp(2)\).  A geometric generator is
\begin{equation}
 Y_5=(S^4,P_1)\times S^1_{\rm R}.
 \label{eq:geometric-generator}
\end{equation}
It is bordant to the mapping torus obtained by gluing the ends of
\(S^4\times[0,1]\) with the nontrivial element of
\(\pi_4(Sp(2))=\ZZ_2\), and also to the \(S^5\) clutching construction.
For a four-dimensional left Weyl fermion in representation \(R\), product
factorization gives
\begin{equation}
 \tau_R(Y_5)=(-1)^{\ind_{S^4}D_R}.
 \label{eq:product-eta-index}
\end{equation}

Restrict Eqs.~\eqref{eq:four-spin4}--\eqref{eq:ten-spin4} to the instanton
\(SU(2)_c\).  The required decompositions and instanton indices are
\begin{align}
 \bm4&\downarrow SU(2)_c=\bm2\oplus\bm1\oplus\bm1,
 &\ind D_{\bm4}&=1,
 \label{eq:index-four}\\
 \bm5&\downarrow SU(2)_c=2\,\bm2\oplus\bm1,
 &\ind D_{\bm5}&=2,
 \label{eq:index-five}\\
 \bm{10}&\downarrow SU(2)_c=\bm3\oplus2\,\bm2\oplus3\,\bm1,
 &\ind D_{\bm{10}}&=4+2=6.
 \label{eq:index-ten}
\end{align}
Here \(2T_{SU(2)}(\bm2)=1\) and
\(2T_{SU(2)}(\bm3)=4\).  Equivalently, in the convention
\(T_{Sp(2)}(\bm4)=1/2\),
\begin{equation}
 2T(\bm4)=1,
 \qquad2T(\bm5)=2,
 \qquad2T(\bm{10})=6.
 \label{eq:dynkin-mod2}
\end{equation}
Therefore
\begin{equation}
 \nu_{\bm4}=1,
 \qquad \nu_{\bm5}=0,
 \qquad \nu_{\bm{10}}=0
 \quad\text{in }\ZZ_2.
 \label{eq:characters}
\end{equation}
The fundamental \(\bm4\) detects the Witten class
\cite{Witten1982,WangWenWitten2019}; the vector \(\bm5\) and adjoint
\(\bm{10}\) do not.

A Dirac \(\bm5\) is two left-Weyl \(\bm5\)'s after charge conjugation.
For two Dirac flavours the exponent is
\begin{equation}
 \nu_{2{\rm D}\,\bm5}
 =2\ {\rm flavours}\times2\ {\rm Weyl/Dirac}\times\nu_{\bm5}
 =4\times0=0\pmod2.
 \label{eq:dirac-character}
\end{equation}
Scalars do not contribute a fermionic Dai--Freed character.

\begin{corollary}[Gauge-bordism gate closes]
The complete dynamical \(Sp(4)\) gauge determinant for the displayed
fermion list evaluates to \(+1\) on the generator of
\(\Omega_5^{\rm Spin}(BSp(2))\).  Thus
\begin{equation}
 \texttt{sp4\_gauge\_bordism\_character\_trivial=true}.
 \label{eq:gauge-gate-closed}
\end{equation}
This conclusion is stronger than merely observing that the fields are
vectorlike: the nonzero bordism group and its representation character have
both been computed.
\end{corollary}

\section{Why the two target mapping-torus phases remain undetermined}

\subsection{Different domains for the two characters}

Put \(X=S^3\times S^2\).  The target-space spin refinement is a character
of
\begin{equation}
 \widetilde\Omega_4^{\rm Spin}(X)=\ZZ_2\oplus\ZZ_2,
 \label{eq:target-bordism}
\end{equation}
with generators
\begin{align}
 G_3&=(S^1_{\rm R}\times S^3,\operatorname{pr}_{S^3},{\rm pt}),
 \label{eq:G3}\\
 G_2&=(T^2_{\rm RR}\times S^2,{\rm pt},\operatorname{pr}_{S^2}).
 \label{eq:G2}
\end{align}
The most general torsion factor is
\begin{equation}
 Z_{\rm tor}(M,g,s)=(-1)^{\epsilon_3\nu_3(M,g)+
                              \epsilon_2\nu_2(M,s)},
 \qquad(\epsilon_3,\epsilon_2)\in\ZZ_2^2.
 \label{eq:target-character}
\end{equation}
Here the invariants have a concrete mod-two definition.  For regular values
\(p_3\in S^3\) and \(p_2\in S^2\), give
\(L_g=g^{-1}(p_3)\) and \(\Sigma_s=s^{-1}(p_2)\) their induced spin
structures and set
\begin{equation}
 \nu_3(M,g)=\dim\ker D_{L_g}\pmod2\in\Omega_1^{\rm Spin}=\ZZ_2,
 \qquad
 \nu_2(M,s)=\operatorname{Arf}(\Sigma_s)\in
 \Omega_2^{\rm Spin}=\ZZ_2.
 \label{eq:target-mod2-definitions}
\end{equation}
Their exponentials are abstract invertible spin-TQFT characters.  They may
be stacked as torsion counterterms/sectors, but Eq.~\eqref{eq:target-character}
does not by itself construct a massive microscopic regulator realizing any
chosen pair.

By contrast, Eq.~\eqref{eq:omega5-result} classifies gauge bundles on
five-manifolds.  There is no canonical map
\begin{equation}
 \widetilde\Omega_4^{\rm Spin}(S^3\times S^2)
 \longrightarrow\Omega_5^{\rm Spin}(BSp(2)).
 \label{eq:no-canonical-map}
\end{equation}
To produce a determinant phase one must specify a regulated mass-family
lift
\begin{equation}
 \mathfrak L:\ (g,s)\longmapsto
 [A(g,s),\Sigma(g,s),\Phi(g,s),\mathcal M(g,s),\{c_a,\Lambda_a,\sigma_a\}]
 \in\mathcal F_{\rm Fred}/\mathcal G,
 \label{eq:mass-family-lift}
\end{equation}
and then pull back the determinant line and its Dai--Freed connection.
The action in Section~2 only contains the elementary
\(A,\Sigma,\Phi\).  It does not specify the composite pion
\(g\)-dependent chiral mass, nor a complete \(s\)-dependent fermion
Yukawa operator, nor the PV continuation along either generator.

\begin{theorem}[Target eta underdetermination]
The computation
\(\Omega_5^{\rm Spin}(BSp(2))=\ZZ_2\) and the trivial gauge character of
the two Dirac \(\bm5\)'s do not determine
\((\epsilon_3,\epsilon_2)\).  With the current field/action data, neither
\(\tau(G_3)\) nor \(\tau(G_2)\) is a defined microscopic number.
\end{theorem}
\begin{proof}
The two putative characters have different bordism domains, as shown by
Eqs.~\eqref{eq:target-bordism} and \eqref{eq:omega5-result}.  A comparison
requires the extra map \(\mathfrak L\) in
Eq.~\eqref{eq:mass-family-lift}.  It is absent.  Independently, the four
characters in Eq.~\eqref{eq:target-character} are distinguished by the two
mod-two indices in Eq.~\eqref{eq:target-mod2-definitions}.  Stacking either
abstract invertible torsion sector has zero de Rham curvature and does not
alter the \(Sp(4)\) gauge character or perturbative anomaly polynomial, but
it flips one target character.  Because the microscopic theory has not
fixed these invertible sectors/counterterms, all four character assignments
remain compatible with the supplied local and gauge-bordism data.  This is
an ambiguity statement, not a construction of four heavy determinants.
\end{proof}

The explicit negative-control table is
\begin{center}
\begin{tabular}{cccc}
\toprule
invertible sector/counterterm & \((\epsilon_3,\epsilon_2)\) & \(Z(G_3)\) & \(Z(G_2)\)\\
\midrule
reference & \((0,0)\) & \(+1\) & \(+1\)\\
\(I_3=(-1)^{\nu_3}\) & \((1,0)\) & \(-1\) & \(+1\)\\
\(I_2=(-1)^{\nu_2}\) & \((0,1)\) & \(+1\) & \(-1\)\\
\(I_3I_2\) & \((1,1)\) & \(-1\) & \(-1\)\\
\bottomrule
\end{tabular}
\end{center}
This table enumerates the character group and proves non-uniqueness of the
present data.  It is neither a menu from which the desired answer may be
selected nor evidence that any row has been generated by the displayed
fermion determinant.

\subsection{Conditional spectral controls}

Define a conditional no-defect reference \({\cal C}_0\) on the retained
branch by
\begin{equation}
 \mathcal K_{{\cal C}_0}(g,s)
 =\slashed D_{A(s)}+
 \mu\bigl(P_Lg+P_Rg^\dagger\bigr),
 \qquad \mu>0,
 \label{eq:conditional-c0}
\end{equation}
with the \(q=-1,0\) branch masses kept constant, a formally positive PV
reference orientation, and no \(I_3\) or \(I_2\) counterterm.  Here \(A(s)\) is the pullback
Hopf connection seen by the charge-one light branch.  This is a formal
\emph{conditional low-energy} parity control; it is not the missing all-field
deep-\(Sp(4)\) lift in Eq.~\eqref{eq:mass-family-lift}.

Two useful parity controls then follow.  In the standard light two-colour
chiral mass family, the \(S^3\) discrete-WZW parity is
\begin{equation}
 N_c\pmod2=2\pmod2=0.
 \label{eq:g3-control}
\end{equation}
For a unit Hopf line over the \(S^2\) factor, the number of light charged
zero-mode copies is
\begin{equation}
 N_cN_f|c_1|=2\cdot2\cdot1=4=0\pmod2.
 \label{eq:g2-control}
\end{equation}
Thus the formal parity count in \({\cal C}_0\) would return \((+1,+1)\).
It is not an evaluated microscopic eta invariant.  Stacking \(I_3\) or
\(I_2\) changes the character without changing the stated
light spectra.  Equations~\eqref{eq:g3-control} and
\eqref{eq:g2-control} therefore cannot be promoted to a unique
microscopic selection.

\begin{gate}[Mapping-torus pair]
Until Eq.~\eqref{eq:mass-family-lift} is supplied and the two APS problems
are evaluated with the same PV convention,
\begin{equation}
 \texttt{target\_mapping\_torus\_eta\_pair\_selected=false}.
 \label{eq:target-pair-false}
\end{equation}
\end{gate}

\section{Heavy thresholds and the fate of
\texorpdfstring{\(k=+2\)}{k=+2}}

\subsection{Weight, multiplicity, mass-sign, and index ledger}

For \(yV>0\), \(m=-yV\), and the positive light regulator mass \(\mu\),
the complete one-flavour branch ledger is
\begin{center}
\begin{tabular}{cccccc}
\toprule
branch & \(U(1)_H\) charge & \(SU(2)_c\) & multiplicity & mass & sign\\
\midrule
light & \(+1\) & \(\bm2\) & 2 & \(\mu\) & \(+\)\\
charged heavy & \(-1\) & \(\bm2\) & 2 & \(-2yV\) & \(-\)\\
neutral heavy & \(0\) & \(\bm1\) & 1 & \(-yV\) & \(-\)\\
\bottomrule
\end{tabular}
\end{center}
There are two identical flavour copies.  Mass signs control theta-like
threshold terms relative to the PV reference, whereas the mixed flavour
coefficient below is a charge trace.  These two ledgers must not be
identified.

For negative real masses and a positive PV reference, the pure unbroken
gauge theta shifts, in units of \(\pi\), are
\begin{align}
 \frac{\Delta\theta_c}{\pi}
 &=N_f\,[2T_{SU(2)}(\bm2)]=2\cdot1=2,
 \label{eq:theta-c}\\
 \frac{\Delta\theta_H}{\pi}
 &=N_f\,[\dim\bm2\,q^2]=2\cdot2\cdot1=4.
 \label{eq:theta-h}
\end{align}
Both exponentiate trivially on integral instanton sectors.  The neutral
branch has no gauge theta shift.  Flipping all PV reference signs changes
these integers by even amounts for the two-flavour spectrum.  This is only
a local pure-gauge theta-periodicity check in the stated normalization.  It
is \emph{not} an APS eta or determinant-ratio matching calculation.  The
unbroken-group Dai--Freed ratio, gravitational phase, spectral flow, and
boundary zero-mode contribution remain open.

As a second negative control, take \(m=+yV\).  Then
\begin{equation}
 (M_+,M_-,M_0)=(2yV,0,yV),
 \label{eq:opposite-tuning}
\end{equation}
so the \(q=-1\), not the \(q=+1\), colour doublet is light.  Its infrared
orientation is \(k=-2\).  This shows that the sign of \(k_{\rm IR}\) is
fixed by the selected branch and generator orientation, not by the
positivity of the full two-flavour measure.

\subsection{Conditional uniform-flavour charge-trace obstruction}

Let \(SU(2)_L\times SU(2)_R\) act uniformly on the two flavour copies of a
complete \(\bm5\).  In the normalization where a retained
\(\bm2_{+1}\) contributes \(+2\), the left mixed coefficient is
\begin{equation}
 \kappa_L^{\rm full}
 =\sum_{w\in\bm5}q_w
 =2(+1)+2(-1)+1(0)=0,
 \qquad \kappa_R^{\rm full}=-\kappa_L^{\rm full}=0.
 \label{eq:full-trace}
\end{equation}
Splitting light and heavy pieces gives
\begin{equation}
 \kappa_L^{\rm light}=+2,
 \qquad
 \kappa_L^{\rm heavy}=-2,
 \qquad
 \kappa_L^{\rm light}+\kappa_L^{\rm heavy}=0.
 \label{eq:threshold-cancel}
\end{equation}
These are representation-theory charge traces, not terms extracted from an
evaluated regulated determinant.  The split acquires threshold meaning only
under the uniform exact-flavour and symmetry-preserving matching hypotheses
stated next.
This equality follows from tracelessness of the embedded semisimple
generator.  It is independent of the signs in
Eq.~\eqref{eq:masses-tuned}.  Indeed, multiplying an \(SU(2)\) chiral mass
map by \(-1\) sends \(g\mapsto-g\), but
\begin{equation}
 (-g)^{-1}\dd(-g)=g^{-1}\dd g,
 \label{eq:minus-g}
\end{equation}
so the normalized \(S^3\) winding form is unchanged.  Mass sign affects
the separate theta ledger, not the charge-trace identity.

\begin{proposition}[Conditional uniform-flavour trace obstruction]
Assume the chiral flavour symmetry acts uniformly on a complete
\(Sp(4)\) \(\bm5\), and the heavy threshold preserves that exact symmetry.
Then the complete-irrep charge trace is zero, while its formal light/heavy
decomposition is \(+2-2\).  Hence any symmetry-preserving matching
construction would have to account for the heavy \(-2\) rather than simply
discard it.  This conditional obstruction does not assert that the actual
heavy determinant has been calculated.
\end{proposition}
\begin{proof}
The assumed uniform flavour assignment makes the relevant representation
trace additive over the five weights.  Equations~\eqref{eq:full-trace} and
\eqref{eq:threshold-cancel} are then exact algebraic identities.  Promoting
them to a statement about determinant phases requires the missing regulated
mass interpolation and APS ratio; no such promotion is made here.
\end{proof}

This is the concrete \(Sp(4)\) realization of the recent semisimple
deep-UV warning: a nonzero mixed WZW/2-group structure requires the
relevant flavour and magnetic symmetries to emerge together below the
semisimple breaking scale, rather than remain exact throughout the deep UV
\cite{DavighiLohitsiri2024}.

The actual operator \(m+y\widehat\Sigma_{\bm5}\) makes this distinction
unavoidable.  A nonzero heavy-branch Dirac mass preserves only the vector
flavour subgroup, not an exact independent
\(SU(2)_L\times SU(2)_R\) acting on the full \(\bm5\).  The candidate can
therefore evade the hypothesis of the proposition only by treating the desired
chiral symmetry as emergent in the tuned light branch.  That is a logically
consistent escape route, but it no longer supplies a deep-UV anomaly or
2-group matching derivation of \(k=+2\); the emergence and threshold
functional must be constructed explicitly.

\subsection{What remains true about physical orientation}

Choose the generator \(H\) so the retained colour doublet has
\(X_q=+1\), and orient the baryon current so a soliton has \(B=+1\).  The
light charged-two-colour sector then has
\begin{equation}
 n_{\rm IR}=N_cX_q=2,
 \qquad
 k_{\rm IR}=n_{\rm IR}B=+2.
 \label{eq:k-ir}
\end{equation}
The opposite branch in Eq.~\eqref{eq:opposite-tuning} gives \(-2\), and
CPT sends \(B\to-B\).  Thus Eq.~\eqref{eq:k-ir} is a physical, declared
infrared orientation.  It is not an all-scale matching theorem.

There are three logically possible repairs:
\begin{enumerate}
\item make \(SU(2)_L\times SU(2)_R\) emergent only after
  \(Sp(4)\to SU(2)_c\times U(1)_H\), while explicitly tracking the
  symmetry-emergence scale and every threshold;
\item add a specified heavy invertible sector whose integer mixed
  functional changes the threshold by \(+2\), then compute its two target
  torsion signs with the same regulator;
\item abandon the uniform complete-irrep flavour assignment and construct
  a different anomaly-consistent UV representation set.
\end{enumerate}
Option 1 is closest to the known deep-UV mechanism, but none has yet been
implemented in this candidate.

\begin{gate}[Full heavy-threshold eta matching]
Although the pure gauge theta shifts in
Eqs.~\eqref{eq:theta-c}--\eqref{eq:theta-h} pass a local periodicity check,
the unbroken-group and gravitational APS determinant ratios, target-dependent
heavy determinant, and all-scale \(+2\) coefficient are not computed.
Therefore
\begin{equation}
 \begin{gathered}
 \texttt{local\_pure\_gauge\_theta\_periodicity\_check=true},\\
 \texttt{unbroken\_group\_aps\_determinant\_ratio\_computed=false},\\
 \texttt{gravitational\_aps\_phase\_computed=false},\\
 \texttt{heavy\_threshold\_eta\_matched=false},\\
 \texttt{k\_plus\_two\_ir\_orientation\_derived=true},\\
 \texttt{k\_plus\_two\_all\_scale\_matched=false}.
 \end{gathered}
 \label{eq:threshold-gates}
\end{equation}
\end{gate}

\section{Mechanical and numerical verification}

The companion script performs exact rational/integer checks and small
spectral controls.  Its independent tests include:
\begin{enumerate}
\item the low-degree \(BSp(2)\) AHSS diagonal and every possible incoming
  and outgoing differential;
\item the \(SU(2)_c\) restriction of \(\bm4,\bm5,\bm{10}\) and mod-two
  Dynkin indices \((1,0,0)\);
\item the four-Weyl multiplicity of two Dirac \(\bm5\)'s;
\item the partial fermionic PV moments and gamma-five controls, together
  with an explicit check that the full regulator gate remains false;
\item both mass tunings, all charge multiplicities, and the local
  pure-gauge theta-periodicity shifts \((2,4)\);
\item the conditional uniform-flavour trace identity \(+2-2=0\), without
  claiming an APS determinant evaluation;
\item finite-mode gamma-five pairing and reference zero-mode parities for
  \(G_3,G_2\);
\item the four abstract target invertible characters, their mod-two-index
  definitions, and all fail-closed booleans.
\end{enumerate}
Passing these checks validates the algebra and the conditional obstruction.  It does
not convert a false physics gate into a true one.

\section{Final theorem and gate ledger}

\begin{theorem}[AP-E6 result]
For the displayed simply connected \(Sp(4)\) field list:
\begin{enumerate}
\item fermionic gamma-five and PV-moment controls are explicit, but the
  complete Euclidean regulator remains open;
\item \(\Omega_5^{\rm Spin}(BSp(2))=\ZZ_2\), and the two Dirac
  \(\bm5\)'s have trivial character on its generator;
\item the target mapping-torus pair remains underdetermined because the
  necessary mass-family lift is absent and abstract invertible torsion
  sectors/counterterms can flip either character;
\item the local pure-gauge theta shifts pass a periodicity check, while a
  uniform complete-irrep flavour hypothesis gives only the conditional
  charge-trace obstruction \(+2-2=0\); no heavy APS ratio is evaluated;
\item the light \(B=+1\) sector has the physical orientation \(+2\), but
  no all-scale \(+2\) matching, radiative protection, or nonperturbative
  \(B=1\) phase is established.
\end{enumerate}
Consequently this work closes the dynamical gauge-bordism subgate, but not
the Euclidean-regulator, heavy-threshold, or full \(Sp(4)\) lane.
\end{theorem}

\begin{longtable}{@{}L{0.42\textwidth}L{0.14\textwidth}L{0.38\textwidth}@{}}
\toprule
Gate & status & evidence\\
\midrule
\endfirsthead
\toprule
Gate & status & evidence\\
\midrule
\endhead
Fermion PV/gamma-five controls & \status{partial} &
Eqs.~\eqref{eq:gamma5-hermiticity}--\eqref{eq:pv-moments}\\
Complete Euclidean \(Sp(4)\) regulator & \status{open} &
mass interpolation, APS boundary/local scheme, and bosonic regulators absent\\
\(\Omega_5^{\rm Spin}(BSp(2))\) & \status{closed} &
AHSS gives the unique surviving \(E^2_{4,1}=\ZZ_2\)\\
Displayed gauge-bordism character & \status{closed} &
\(\nu_{\bm5}=0\), hence two Dirac \(\bm5\)'s are trivial\\
Target \((\epsilon_3,\epsilon_2)\) & \status{open} &
mass-family lift absent; four invertible characters remain possible\\
Local pure-gauge theta periodicity & \status{checked} &
local shifts \((2\pi,4\pi)\) exponentiate trivially\\
Unbroken/gravitational APS threshold & \status{open} &
no regulated determinant ratio, spectral flow, or boundary analysis\\
Target-dependent heavy threshold & \status{open} &
uniform exact flavour only conditionally gives the trace \(+2-2=0\)\\
All-scale \(k=+2\) & \status{open} &
only the light \(B=+1\) orientation is derived\\
Radiative protection & \status{open} &
\(m=-yV\) remains an unprotected cancellation\\
Strong charged \(B=1\) phase & \status{open} &
no dynamical lattice/continuum soliton proof in this note\\
Full lane / Route-E portal & \status{blocked} &
target eta, threshold, naturalness, and dynamics gates remain false\\
\bottomrule
\end{longtable}

The machine-readable status is
\begin{equation}
 \begin{gathered}
 \texttt{fermion\_pv\_gamma5\_controls\_verified=true},\\
 \texttt{sp4\_euclidean\_regulator\_complete=false},\\
 \texttt{omega5\_spin\_bsp4\_computed=true},\\
 \texttt{target\_mapping\_torus\_eta\_pair\_selected=false},\\
 \texttt{local\_pure\_gauge\_theta\_periodicity\_check=true},\\
 \texttt{heavy\_threshold\_eta\_matched=false},\\
 \texttt{threshold\_radiatively\_protected=false},\\
 \texttt{strong\_b1\_phase\_proven=false},\\
 \texttt{lane\_closed=false},\qquad
 \texttt{physics\_promotion\_allowed=false}.
 \end{gathered}
 \label{eq:machine-status}
\end{equation}

\section{Next discriminating calculations}

The following alternatives are ordered by information gain.
\begin{enumerate}
\item \textbf{Emergent-flavour construction.}  Specify the fields that
  split the full \(\bm5\) flavour symmetry at the \(Sp(4)\) breaking scale,
  derive the exact light flavour symmetry, and calculate the induced
  integer threshold functional.  This directly tests whether the
  semisimple no-go assumptions are escaped rather than merely violated.
\item \textbf{Common target mass map.}  Construct
  \(\mathfrak L(g,s)\) in Eq.~\eqref{eq:mass-family-lift} on both
  \(G_3\) and \(G_2\), including every heavy branch and PV sign.  A single
  lattice spectral-flow implementation should then compute both eta
  phases without independent convention changes.
\item \textbf{Alternative representation set.}  Search for complete
  \(Sp(4)\) representations plus spectators whose exact flavour-weighted
  charge trace is \(+2\), whose \(\ZZ_2\) gauge-bordism character vanishes,
  and whose scalar sector still realizes the charge-one/exactly-two
  operator.  The trace and mod-two-index constraints can be imposed before
  any dynamical scan.
\end{enumerate}
No degree-one portal should be built until at least one full lane, including
its nonperturbative \(B=1\) gate, closes.

\bibliographystyle{unsrt}
\bibliography{ap_e6_sp4_eta_bordism_threshold}

\end{document}
